\section{Independent PRORL Benchmark and Generalization Assessment}
\label{sec:prorl-benchmark}
We evaluated the published PRORL policy in its upstream four-node synthetic
environment as a separate reproducibility benchmark. This environment differs
from the packet-service elasticity simulator used in the earlier manuscript.
Consequently, its reward values are not compared numerically with those of the
proposed elasticity controller, and this experiment does not constitute a
matched comparison between the two methods.

The training-seed-10 run completed 89,400 loop steps, including 3,000 bootstrap
steps. Model selection used validation seed 100. The saved selected checkpoint
was then evaluated using seeds 1000, 1100, and 1200, each for 168 steps. An
additional evaluation used seeds 2000--2009 without further training or model
selection. The source model and results were checked against the archived
SHA-256 fingerprints. Source revision, dependency versions, configuration and
evaluation outputs are retained in the accompanying repository.

For each policy, the resource initialization, emulator configuration, reward
and evaluation length were matched. Only the agent configuration changed.
Resource inputs were reconstructed from the original YAML configuration to
avoid applying the upstream derived-allocation calculation twice to an
already-mutated configuration export. The selected PRORL state was explicitly
loaded into the evaluation agent. These checks establish the stated software
comparison; they do not establish physical-model calibration.

\begin{table}[t]
\centering
\caption{Additional evaluation at the original demand amplitude (ten evaluation
seeds; one trained checkpoint). Metrics retain their upstream definitions.}
\label{tab:prorl-original}
\begin{tabular}{lrrr}
\hline
Metric & PRORL & Greedy & Heuristic\\
\hline
Utility & 137.120 & 116.670 & 117.009\\
Surplus & 301.7 & 155.6 & 157.9\\
Cost & 19.4 & 235.8 & 236.1\\
Hours satisfied & 149.4 & 115.2 & 116.4\\
Utilization & 0.9100 & 0.9492 & 0.9486\\
\hline
\end{tabular}
\end{table}

At the training demand amplitude of 80, PRORL achieved greater utility than
both comparators on all ten additional evaluation seeds. The mean paired
utility differences were 20.450 against Greedy and 20.111 against Heuristic.
However, surplus increased and utilization decreased
(Table~\ref{tab:prorl-original}). Thus, the observation supports a
reward-specific advantage in this configuration rather than uniform
dominance across resource metrics.

We next changed only the synthetic demand amplitude to 60 and 100, retaining
the saved model and the same ten evaluation seeds. The amplitudes were fixed
before these runs. Because these seeds had already been examined in the
same-load experiment, this assessment is an exploratory paired stress test,
rather than an untouched confirmatory test set. No policy was tuned using the
stress-test outcomes.

\begin{table}[t]
\centering
\caption{Utility under changed synthetic demand amplitudes. No retraining was
performed.}
\label{tab:prorl-shift}
\begin{tabular}{rrrr}
\hline
Amplitude & PRORL & Greedy & Heuristic\\
\hline
60 & 106.000 & 121.676 & 121.677\\
100 & 108.964 & 110.844 & 111.800\\
\hline
\end{tabular}
\end{table}

The original utility advantage did not persist at either changed amplitude
(Table~\ref{tab:prorl-shift}). At amplitude 60, PRORL's mean surplus reached
925.5, compared with 160.7 for Greedy, while utilization was 0.7234 versus
0.9329. At amplitude 100, the corresponding utilization values were 0.9974
and 0.9992, and PRORL's utility remained lower. These outcomes identify a
generalization limitation of this saved checkpoint; they do not isolate its
cause. Establishing a mechanism requires inspection of reward terms and
action histories, with matched ablations.

The runtime was restored using Python 3.11.17, with all recorded dependency
versions matching the saved environment list. The original interpreter patch
version was not established. This qualification is retained with the results.
The metrics are synthetic service and resource counters; cost is not measured
energy, and the experiment does not measure packet latency or validate
software-defined networking interfaces.

Only one completed training seed supports the reported tables. An independent
seed-20 replication is in progress and is excluded from the numerical results
until training, final evaluation and integrity checks finish. Additional
evaluation seeds cannot substitute for independent training replications.
The benchmark establishes neither a new allocation algorithm nor the novelty
of the thesis contribution. A direct PRORL comparison in the manuscript's
service model requires a documented mapping of observations, actions,
resource budgets and objectives, followed by training and evaluation under
that shared model.
